% =====================================================================
% pertemuan-02.tex
% Pertemuan 2 - Teori Bahasa Formal dan Hierarki Chomsky
% Sub-CPMK 1.4 - Teori terapan + kuis konsep
% =====================================================================
\documentclass[aspectratio=169,10pt,t]{beamer}
\usepackage{config/beamerthemeubb}

\UBBmeeting{2}

\begin{document}

\UBBtitleframe{2}{Teori Bahasa Formal dan Hierarki Chomsky}{Teori terapan \textperiodcentered{} Kuis konsep}

% ---------------------------------------------------------------------
\begin{frame}{Capaian Pertemuan}
  \textbf{CPMK induk:} CPMK 01 --- Menjelaskan tahapan kompilasi
  (\textit{lexical}, \textit{syntax}, \textit{semantic analysis}).

  \vspace{4pt}
  \begin{block}{Sub-CPMK yang dinilai pada pertemuan ini}
    \small
    \textbf{Sub-CPMK 1.4:} Membuat tabel pemetaan hierarki Chomsky
    (Tipe 3--0) ke jenis automata dan menguji klasifikasi lima contoh
    grammar ke kelas yang tepat.
  \end{block}

  \vspace{4pt}
  \begin{exampleblock}{Luaran yang diperiksa}
    \footnotesize
    Tabel empat baris yang lengkap, klasifikasi lima grammar beserta
    alasan berbasis bentuk aturan, dan hasil uji string anggota serta
    bukan anggota.
  \end{exampleblock}
\end{frame}

% ---------------------------------------------------------------------
\begin{frame}{Agenda Pertemuan (150 menit)}
  \small
  \begin{tabular}{@{}p{1.4cm}p{12.3cm}@{}}
    \toprule
    \textbf{Menit} & \textbf{Kegiatan} \\
    \midrule
    10 & Pembuka dan orientasi capaian \\
    35 & Alfabet, string, bahasa, serta tata bahasa dan derivasi \\
    40 & Hierarki Chomsky Tipe 3--0 dan automata pengenalnya \\
    25 & Tabel pemetaan dan klasifikasi lima grammar \\
    20 & Kuis konsep dan pembahasan \\
    20 & Diskusi implikasi bagi desain lexer dan parser \\
    \bottomrule
  \end{tabular}

  \vspace{6pt}
  {\scriptsize\color{ubbblue!75}Pertemuan teori terapan: ceramah singkat,
  lokakarya tabel, lalu kuis konsep.}
\end{frame}

% ---------------------------------------------------------------------
\begin{frame}{Posisi Pertemuan pada Rantai Proyek}
  \UBBchain{2}

  \vspace{2pt}
  \begin{itemize}
    \item Pertemuan ini menyediakan landasan formal untuk simpul
          \textit{grammar} pada rantai proyek.
    \item Teori bahasa menjelaskan mengapa lexer cukup memakai otomata
          hingga dan parser memerlukan otomata pushdown.
    \item Pertemuan ini tidak menambah increment kode; keluarannya
          berupa tabel dan uji klasifikasi.
  \end{itemize}
\end{frame}

% ---------------------------------------------------------------------
\begin{frame}{Alfabet, String, Bahasa, dan Operasinya}
  \begin{block}{Definisi dasar}
    \small
    Alfabet \(\Sigma\) adalah himpunan berhingga dan tidak kosong dari
    simbol. String adalah barisan berhingga simbol; panjangnya
    \(|w|\), dan string kosong \(\varepsilon\) panjangnya nol.
    Himpunan seluruh string adalah \(\Sigma^{*}\). Bahasa
    \(L \subseteq \Sigma^{*}\) adalah himpunan string.
  \end{block}

  \vspace{2pt}
  \begin{columns}[T,onlytextwidth]
    \begin{column}{0.50\textwidth}
      \footnotesize
      \textbf{Operasi bahasa}
      \begin{itemize}
        \item gabungan: \(L_1 \cup L_2\);
        \item konkatenasi: \(L_1L_2 = \{xy \mid x \in L_1, y \in L_2\}\);
        \item bintang Kleene: \(L^{*} = \bigcup_{k \ge 0} L^k\);
        \item \textit{positive closure}: \(L^{+} = \bigcup_{k \ge 1} L^k\).
      \end{itemize}
    \end{column}
    \begin{column}{0.46\textwidth}
      \begin{exampleblock}{Contoh}
        \footnotesize
        \(\Sigma=\{a,b\}\), \(L=\{a,bb\}\).\\
        \(L^{*}\) memuat \(\varepsilon, a, bb, aa, abb, bba, bbbb,
        \ldots\); \(L^{+}\) sama tanpa \(\varepsilon\).
      \end{exampleblock}
    \end{column}
  \end{columns}
\end{frame}

% ---------------------------------------------------------------------
% ---------------------------------------------------------------------
\begin{frame}{Tata Bahasa sebagai Empat-Tupel}
  \begin{block}{Definisi}
    \(G = (V, T, P, S)\), dengan \(V\) himpunan nonterminal, \(T\)
    himpunan terminal, \(P\) himpunan aturan produksi, dan \(S \in V\)
    simbol awal. Himpunan \(V\) dan \(T\) saling lepas.
  \end{block}

  \begin{itemize}
    \item Aturan berbentuk \(\alpha \to \beta\), dengan \(\alpha\)
          memuat paling sedikit satu nonterminal.
    \item Jika \(\alpha = A \in V\), aturan disebut bebas konteks.
    \item Bahasa yang dibangkitkan:
          \(L(G) = \{w \in T^{*} \mid S \Rightarrow^{*} w\}\).
    \item Dua tata bahasa ekivalen jika \(L(G_1)=L(G_2)\).
  \end{itemize}

  \vspace{2pt}
  {\footnotesize\color{ubbblue!80}Contoh:
  \(S \to AB\), \(A \to aA \mid a\), \(B \to bB \mid b\).}
\end{frame}

% ---------------------------------------------------------------------
\begin{frame}{Derivasi dan Bentuk Kalimat}
  \begin{itemize}
    \item Satu langkah: \(\alpha A\beta \Rightarrow \alpha\gamma\beta\)
          jika \(A \to \gamma \in P\).
    \item \(\Rightarrow^{*}\) berarti nol atau lebih langkah;
          \(\Rightarrow^{+}\) berarti satu atau lebih langkah.
    \item Bentuk kalimat adalah string yang dapat diturunkan dari \(S\).
  \end{itemize}

  \vspace{3pt}
  \begin{exampleblock}{Contoh: string \(aabb\) dari
    \(S \to AB,\ A \to aA \mid a,\ B \to bB \mid b\)}
    \footnotesize
    \textbf{Leftmost:}
    \(S \Rightarrow AB \Rightarrow aAB \Rightarrow aaB
     \Rightarrow aabB \Rightarrow aabb\).\\
    \textbf{Rightmost:}
    \(S \Rightarrow AB \Rightarrow AbB \Rightarrow Abb
     \Rightarrow aAbb \Rightarrow aabb\).
  \end{exampleblock}

  \vspace{2pt}
  {\footnotesize Keduanya menghasilkan pohon urai yang sama karena
  tata bahasa ini tidak ambigu.}
\end{frame}

% ---------------------------------------------------------------------
\begin{frame}{Pohon Urai}
  \begin{columns}[T,onlytextwidth]
    \begin{column}{0.48\textwidth}
      \centering
      \begin{tikzpicture}[
          nd/.style={draw=ubbblue!40,fill=ubbblue!7,rounded corners=2pt,
            minimum width=0.55cm,minimum height=0.48cm,font=\small,
            inner sep=1pt},
          ar/.style={thick,draw=ubbblue!55}
        ]
        \node[nd] (S) at (0,0) {\(S\)};
        \node[nd] (A) at (-1.5,-0.9) {\(A\)};
        \node[nd] (B) at (1.5,-0.9) {\(B\)};
        \node[nd] (a1) at (-2.4,-1.8) {\(a\)};
        \node[nd] (A2) at (-0.9,-1.8) {\(A\)};
        \node[nd] (b1) at (0.9,-1.8) {\(b\)};
        \node[nd] (B2) at (2.4,-1.8) {\(B\)};
        \node[nd] (a2) at (-0.9,-2.7) {\(a\)};
        \node[nd] (b2) at (2.4,-2.7) {\(b\)};
        \draw[ar] (S) -- (A);  \draw[ar] (S) -- (B);
        \draw[ar] (A) -- (a1); \draw[ar] (A) -- (A2);
        \draw[ar] (B) -- (b1); \draw[ar] (B) -- (B2);
        \draw[ar] (A2) -- (a2);
        \draw[ar] (B2) -- (b2);
      \end{tikzpicture}
    \end{column}
    \begin{column}{0.48\textwidth}
      \begin{itemize}
        \item Akar: simbol awal \(S\).
        \item Simpul dalam: nonterminal.
        \item Daun: terminal; dibaca kiri ke kanan menghasilkan
              \(aabb\).
        \item Satu string dapat memiliki dua pohon berbeda jika tata
              bahasa ambigu; hal ini dibahas pada Pertemuan 7.
      \end{itemize}
    \end{column}
  \end{columns}
\end{frame}

% ---------------------------------------------------------------------
\begin{frame}{Hierarki Chomsky: Empat Kelas Bersarang}
  \begin{center}
  \begin{tikzpicture}
    \fill[ubbblue!4,rounded corners=4pt] (0,0) rectangle (12.4,5.0);
    \draw[ubbblue!35,thick,rounded corners=4pt] (0,0) rectangle (12.4,5.0);
    \node[anchor=north west,font=\scriptsize\bfseries,text=ubbblue]
      at (0.18,4.82) {Tipe 0: rekursif terbilang --- mesin Turing};

    \fill[ubbblue!9,rounded corners=4pt] (0.5,0.5) rectangle (11.9,4.1);
    \draw[ubbblue!40,thick,rounded corners=4pt] (0.5,0.5) rectangle (11.9,4.1);
    \node[anchor=north west,font=\scriptsize\bfseries,text=ubbblue]
      at (0.68,3.92) {Tipe 1: sensitif konteks --- otomata terbatas linear};

    \fill[ubbblue!15,rounded corners=4pt] (1.0,1.0) rectangle (11.4,3.2);
    \draw[ubbblue!45,thick,rounded corners=4pt] (1.0,1.0) rectangle (11.4,3.2);
    \node[anchor=north west,font=\scriptsize\bfseries,text=ubbblue]
      at (1.18,3.02) {Tipe 2: bebas konteks --- otomata pushdown};

    \fill[ubbblue!25,rounded corners=4pt] (1.5,1.5) rectangle (10.9,2.3);
    \draw[ubbblue!55,thick,rounded corners=4pt] (1.5,1.5) rectangle (10.9,2.3);
    \node[anchor=north west,font=\scriptsize\bfseries,text=ubbblue]
      at (1.68,2.12) {Tipe 3: regular --- otomata hingga};
  \end{tikzpicture}
  \end{center}

  \vspace{1pt}
  {\footnotesize Semakin longgar bentuk aturan, semakin besar daya
  ekspresif kelas bahasa; setiap kelas termuat dalam kelas di atasnya.}
\end{frame}

% ---------------------------------------------------------------------
\begin{frame}{Tipe 3: Tata Bahasa Regular}
  \begin{block}{Bentuk aturan}
    \small
    \(A \to aB\), \(A \to a\), atau \(A \to \varepsilon\), dengan
    \(A,B \in V\) dan \(a \in T\) (\textit{right-linear}). Bentuk
    \textit{left-linear} \(A \to Ba\) menghasilkan kelas bahasa yang
    sama.
  \end{block}

  \begin{itemize}
    \item Ekivalen dengan ekspresi reguler dan otomata hingga.
    \item Cocok untuk token: identifier, bilangan, operator.
    \item Contoh bahasa: \(a^{*}b^{*}\) dan himpunan identifier.
  \end{itemize}

  \vspace{2pt}
  \begin{exampleblock}{Identifier: huruf/garis bawah, lalu huruf/angka}
    \footnotesize
    \(I \to hI \mid \_I \mid h\), \quad
    \(R \to hR \mid dR \mid \varepsilon\),\\
    dengan \(h\) huruf, \(d\) angka, dan \(R\) sisa identifier.
  \end{exampleblock}
\end{frame}

% ---------------------------------------------------------------------
\begin{frame}{Tipe 2: Tata Bahasa Bebas Konteks}
  \begin{block}{Bentuk aturan}
    \small
    \(A \to \alpha\), dengan \(A \in V\) dan
    \(\alpha \in (V \cup T)^{*}\). Kelas ini dikenali oleh otomata
    pushdown (\textit{pushdown automaton}, PDA).
  \end{block}

  \begin{itemize}
    \item Mampu menangani struktur bersarang: kurung, blok, pemanggilan.
    \item Contoh: \(S \to aSb \mid \varepsilon\) membangkitkan
          \(\{a^n b^n \mid n \ge 0\}\).
    \item Contoh ekspresi:
          \(E \to E+T \mid T\), \(T \to T*F \mid F\),
          \(F \to (E) \mid \code{id}\).
    \item Tertutup terhadap gabungan, konkatenasi, dan bintang Kleene;
          tidak tertutup terhadap irisan dan komplemen.
  \end{itemize}
\end{frame}

% ---------------------------------------------------------------------
\begin{frame}{Tipe 1: Tata Bahasa Sensitif Konteks}
  \begin{block}{Bentuk aturan}
    \small
    \(\alpha A\beta \to \alpha\gamma\beta\), dengan
    \(\gamma \neq \varepsilon\). Nonterminal \(A\) diganti hanya dalam
    konteks \(\alpha\) dan \(\beta\), dan panjang string tidak boleh
    berkurang.
  \end{block}

  \begin{itemize}
    \item Dikenali oleh otomata terbatas linear
          (\eng{linear bounded automaton}).
    \item Contoh: \(\{a^n b^n c^n \mid n \ge 1\}\), yang memerlukan
          pencocokan tiga kelompok berjumlah sama.
    \item Satu stack tidak cukup; diperlukan batas pita sepanjang
          masukan.
  \end{itemize}
\end{frame}

% ---------------------------------------------------------------------
\begin{frame}{Tipe 0: Tata Bahasa Tanpa Batas}
  \begin{block}{Bentuk aturan}
    \small
    \(\alpha \to \beta\) tanpa batasan bentuk, asalkan \(\alpha\)
    memuat paling sedikit satu nonterminal.
  \end{block}

  \begin{itemize}
    \item Ekivalen dengan mesin Turing.
    \item Mencakup bahasa yang dapat dikenali secara rekursif
          terbilang.
    \item Sebagian bahasa pada kelas ini tidak dapat diputuskan.
    \item Contoh klasik: \eng{halting problem} --- tidak ada algoritma
          yang menentukan apakah program sembarang berhenti untuk semua
          masukan.
  \end{itemize}
\end{frame}

% ---------------------------------------------------------------------
\begin{frame}{Tabel Pemetaan Hierarki Chomsky (Artefak Sub-CPMK 1.4)}
  \small
  \setlength{\tabcolsep}{4pt}
  \begin{tabular}{@{}p{0.8cm}p{2.7cm}p{4.6cm}p{3.3cm}p{1.9cm}@{}}
    \toprule
    \textbf{Tipe} & \textbf{Kelas bahasa} & \textbf{Bentuk aturan} &
    \textbf{Automata} & \textbf{Contoh} \\
    \midrule
    3 & Regular & \(A \to aB\), \(A \to a\) & DFA/NFA &
      \(a^{*}b^{*}\) \\
    2 & Bebas konteks & \(A \to \alpha\) & PDA &
      \(\{a^n b^n\}\) \\
    1 & Sensitif konteks & \(\alpha A\beta \to \alpha\gamma\beta\) &
      Otomata terbatas linear & \(\{a^n b^n c^n\}\) \\
    0 & Rekursif terbilang & \(\alpha \to \beta\) & Mesin Turing &
      bahasa yang dikenali mesin Turing \\
    \bottomrule
  \end{tabular}

  \vspace{6pt}
  {\footnotesize\color{ubbblue!80}Setiap baris harus konsisten: contoh
  bahasa pada kolom terakhir benar-benar dikenali oleh automata pada
  kolom keempat. Tabel mengikuti Guntoro Bab 6; nama automata
  diperjelas dari Thain 4.6 bila perlu. Tabel inilah artefak
  Sub-CPMK 1.4.}
\end{frame}

% ---------------------------------------------------------------------
\begin{frame}{Contoh Berpenyelesaian: Klasifikasi Lima Grammar}
  \scriptsize
  \setlength{\tabcolsep}{4pt}
  \begin{tabular}{@{}p{0.5cm}p{5.4cm}p{1.4cm}p{6.4cm}@{}}
    \toprule
    & \textbf{Aturan produksi} & \textbf{Tipe} & \textbf{Alasan} \\
    \midrule
    (a) & \(S \to aS \mid bS \mid \varepsilon\) & 3 &
      Semua aturan berbentuk \(A \to aB\) atau \(A \to \varepsilon\). \\
    (b) & \(S \to aSb \mid \varepsilon\) & 2 &
      Aturan \(A \to \alpha\); membangkitkan \(\{a^n b^n\}\). \\
    (c) & \(S \to aSBC \mid aBC\), \(CB \to BC\), \(aB \to ab\),
      \(bB \to bb\), \(bC \to bc\), \(cC \to cc\) & 1 &
      Aturan \(CB \to BC\) bergantung konteks; panjang tidak berkurang;
      membangkitkan \(\{a^n b^n c^n\}\). \\
    (d) & \(S \to AB\), \(AB \to BA\), \(A \to a\), \(B \to b\) & 0 &
      \(AB \to BA\) tidak memenuhi batasan sensitif konteks. \\
    (e) & \(S \to Sa \mid b\) & 3 &
      Aturan \(A \to Ba\), yaitu \textit{left-linear}. \\
    \bottomrule
  \end{tabular}

  \vspace{5pt}
  {\scriptsize\color{ubbblue!75}Uji tambahan: setiap grammar diuji
  dengan dua string anggota dan satu string bukan anggota.}
\end{frame}

% ---------------------------------------------------------------------
\begin{frame}{Implikasi bagi Desain Lexer dan Parser}
  \small
  \setlength{\tabcolsep}{4pt}
  \begin{tabular}{@{}p{3.0cm}p{3.5cm}p{7.0cm}@{}}
    \toprule
    \textbf{Fase} & \textbf{Formalisme} & \textbf{Alasan} \\
    \midrule
    Leksikal & Regular/DFA & Struktur token sederhana dan efisien \\
    Sintaksis & CFG/PDA & Struktur bersarang memerlukan stack \\
    Semantik & \textit{Attribute grammar} & Memerlukan konteks dan tipe \\
    Validasi & Aturan dan traversal & Memerlukan analisis lintas simpul \\
    \bottomrule
  \end{tabular}

  \vspace{5pt}
  \begin{block}{Peringatan desain}
    \footnotesize
    Jangan memakai ekspresi reguler untuk memeriksa struktur bersarang
    seperti kurung seimbang. Ekspresi reguler tidak memiliki memori tak
    terbatas; gunakan parser atau PDA.
  \end{block}
\end{frame}

% ---------------------------------------------------------------------
\begin{frame}{Kuis Konsep}
  \small
  \begin{enumerate}
    \setlength{\itemsep}{3pt}
    \item Apa perbedaan string dan bahasa?
    \item Tata bahasa \(S \to aS \mid bS \mid \varepsilon\) termasuk
          Tipe berapa? Sebutkan alasannya.
    \item Mengapa \(\{a^n b^n \mid n \ge 0\}\) bukan bahasa regular?
          Mesin apa yang mengenalinya?
    \item Automata apa yang mengenali bahasa sensitif konteks?
    \item Mengapa lexer cukup memakai Tipe 3, sedangkan parser
          memerlukan Tipe 2?
  \end{enumerate}

  \vspace{4pt}
  {\footnotesize\color{ubbblue!75}Kerjakan mandiri selama 8 menit,
  lalu bahas bersama. Kunci tersedia pada versi handout.}
\end{frame}

\mode<handout>{
\begin{frame}{Kunci Kuis Konsep}
  \footnotesize
  \begin{enumerate}
    \setlength{\itemsep}{2pt}
    \item String adalah barisan simbol; bahasa adalah himpunan string
          atas suatu alfabet.
    \item Tipe 3 (regular), karena setiap aturan berbentuk
          \(A \to aB\) atau \(A \to \varepsilon\).
    \item Regular tidak dapat menghitung jumlah \(a\) dan \(b\) yang
          tak terbatas; bahasa ini dikenali otomata pushdown.
    \item Otomata terbatas linear (\eng{linear bounded automaton}).
    \item Token memiliki struktur sederhana sehingga cukup DFA;
          struktur program bersarang memerlukan stack PDA.
  \end{enumerate}
\end{frame}
}

% ---------------------------------------------------------------------
\begin{frame}{Kegiatan Mahasiswa di Kelas}
  \begin{block}{Lokakarya: tabel pemetaan dan uji lima grammar}
    \footnotesize
    \begin{enumerate}
      \setlength{\itemsep}{1pt}
      \item Lengkapi tabel Tipe 3--0: kelas bahasa, bentuk aturan,
            automata, dan contoh bahasa.
      \item Tentukan tipe lima grammar yang diberikan beserta alasannya.
      \item Uji setiap grammar dengan minimal tiga string: dua anggota
            dan satu bukan anggota.
      \item Cocokkan hasil uji dengan automata yang dipetakan.
      \item Kumpulkan tabel, klasifikasi, dan hasil uji.
    \end{enumerate}
  \end{block}

  \vspace{4pt}
  {\footnotesize\color{ubbblue!80}Bekerja berpasangan; satu pasangan
  memeriksa tabel pasangan lain sebelum dikumpulkan.}
\end{frame}

% ---------------------------------------------------------------------
\begin{frame}{Tugas dan Artefak}
  \begin{block}{Artefak Sub-CPMK 1.4}
    Tabel pemetaan hierarki Chomsky Tipe 3--0 ke jenis automata,
    disertai klasifikasi lima contoh grammar dan hasil uji string.
  \end{block}

  \small
  \begin{itemize}
    \item Tabel empat baris lengkap: tipe, kelas bahasa, bentuk aturan,
          automata, dan contoh bahasa.
    \item Lima grammar terklasifikasi dengan alasan berbasis bentuk
          aturan.
    \item Setiap grammar diuji dengan dua string anggota dan satu string
          bukan anggota.
    \item Tidak ada contoh bahasa yang salah kelas.
  \end{itemize}

  \vspace{3pt}
  {\footnotesize\color{ubbblue!75}Artefak ini menjadi bukti Sub-CPMK 1.4
  pada Kuis Konsep (3\%) dan UTS Praktik (14\%).}
\end{frame}

% ---------------------------------------------------------------------
\begin{frame}{Rangkuman dan Jembatan}
  \begin{itemize}
    \item Alfabet, string, dan bahasa didefinisikan secara formal;
          bahasa adalah himpunan string atas alfabet.
    \item Tata bahasa \(G=(V,T,P,S)\) membangkitkan bahasa melalui
          derivasi; leftmost dan rightmost dapat menghasilkan pohon
          yang sama.
    \item Hierarki Chomsky bersarang: Tipe 3 regular, Tipe 2 bebas
          konteks, Tipe 1 sensitif konteks, Tipe 0 tanpa batas.
    \item Setiap tipe bersesuaian dengan automata: DFA/NFA, PDA,
          otomata terbatas linear, dan mesin Turing.
    \item Pilihan formalisme mengikuti struktur masalah: token memakai
          Tipe 3, sintaksis memakai Tipe 2, semantik melampaui Tipe 2.
  \end{itemize}

  \vspace{4pt}
  \begin{exampleblock}{Pertemuan berikutnya}
    Pertemuan 3 memulai analisis leksikal: token, lexeme, pattern, dan
    ekspresi reguler. Teori Tipe 3 berubah menjadi tabel ekspresi
    reguler untuk setiap kategori token subset bahasa C.
  \end{exampleblock}
\end{frame}

\end{document}
